\documentclass[aspectratio=169,11pt]{beamer}
\usepackage[UTF8,fontset=none]{ctex}
\setCJKmainfont{Microsoft JhengHei}
\setCJKsansfont{Microsoft JhengHei}
\usepackage{amsmath,amssymb,booktabs,array}
\usepackage{pgfplots}
\pgfplotsset{compat=1.18}
\definecolor{navy}{HTML}{183750}
\definecolor{teal}{HTML}{007E87}
\definecolor{orange}{HTML}{B45A20}
\setbeamercolor{normal text}{fg=navy,bg=white}
\setbeamercolor{structure}{fg=teal}
\setbeamercolor{frametitle}{fg=navy}
\setbeamercolor{alerted text}{fg=orange}
\setbeamerfont{frametitle}{size=\Large,series=\bfseries}
\setbeamertemplate{navigation symbols}{}
\setbeamertemplate{itemize items}[circle]
\setbeamersize{text margin left=8mm,text margin right=8mm}
\setbeamertemplate{footline}{\hspace{8mm}\scriptsize Handout 2 自讀筆記\hfill\insertsectionhead\hfill\insertframenumber/\inserttotalframenumber\hspace{8mm}\vspace{3mm}}
\newcommand{\src}[1]{\par\vfill{\scriptsize\color{teal}#1}}
\newcommand{\slide}[3]{\begin{frame}[t]{#1}#3\src{#2}\end{frame}}
\newcommand{\bi}{\begin{itemize}\setlength{\itemsep}{5pt}}
\newcommand{\ei}{\end{itemize}}
\newcommand{\key}[1]{\par\medskip{\color{teal}\bfseries #1}\par\smallskip}
\newcommand{\dd}{\mathrm{d}}
\newcommand{\bookslide}[3]{\slide{#1}{課本：#2。頁碼為書本印刷頁碼}{#3}}
% Original images extracted losslessly from the supplied textbook PDF.
\newcommand{\originalfigure}[6]{
\begin{frame}[t]{課本原圖 #1：#2}
\centering
\includegraphics[width=.94\linewidth,height=5.05cm,keepaspectratio]{figures/#4}\par
\vspace{2mm}\raggedright\small #5\par
\vfill{\scriptsize\color{teal}原圖：Collinson, 第 3 版，Fig. #1，p. #3。#6}\par
\end{frame}}


\newcommand{\handoutfigure}[3]{
\begin{frame}[t]{講義原頁 #1：#2}
\centering
\includegraphics[width=.97\linewidth,height=5.65cm,keepaspectratio]{figures/#3}\par
\vfill{\scriptsize\color{teal}原圖與文字：賴盈誌老師提供之 Handout，PDF 第 #1 頁。}\par
\end{frame}}
\title{航電系統架構與數位整合}
\subtitle{Handout 2\\繁體中文自讀詳解}
\author{依賴盈誌老師課程講義整理}
\date{CA7528 Digital Avionics\quad Fall 2026}
\begin{document}
\begin{frame}\titlepage\centering\small 概念、原圖、推導、算例與自我檢查\end{frame}
\section{閱讀地圖}
\slide{範圍與閱讀方式}{Handout 2，PDF pp. 1--57}{
\bi
\item 類比至數位、資料匯流排、CNS、IMA、BIT 與維修。
\item 原頁碼一律指講義 PDF 頁序。標示「課本」的頁碼則是 Collinson 第三章印刷頁碼。
\item 先讀概念，再對照原圖。自我檢查先遮住參考解答，再獨立計算或口頭說明。
\item 自編算例與補充推導依題內假設使用。歷史機型、服務與表格依原教材背景理解。
\ei
}
\slide{分次閱讀安排}{Handout 2 閱讀地圖}{
\small
\bi
\item \textbf{航電與演進}：航電用途與工程單位、外界目視與儀表資訊。
\item \textbf{獨立類比架構}：獨立儀表與類比傳輸。
\item \textbf{駕駛艙與分工}：類比與數位顯示的操作差異。
\item \textbf{資料匯流排}：數位整合與共享資料、控制層級與更新率。
\item \textbf{系統功能}：性能、可用性與機地分工。
\item \textbf{CNS/ATM}：傳統與衛星資料鏈 CNS。
\ei
}
\slide{分次閱讀安排（續）}{Handout 2 閱讀地圖}{
\small
\bi
\item \textbf{模組化與顯示}：IMA、LRU 與微電子技術、不同機型的座艙閱讀方法。
\item \textbf{維修與 BIT}：間歇故障、紀錄與更換。
\item \textbf{整合複習}：端到端航電分析。
\ei
}
\section{航電與演進}
\slide{航電用途與工程單位}{Handout 2，PDF pp. 1--6；概念整理與補充推導}{
\small
\bi
\item 講義由 avionics 的定義與設備實例開始，再比較不同年代的儀表及整合程度。
\item 歷史成本比例與飛機性能數字代表講義案例，不能直接當作所有機型的固定比例。
\item 一海里為 $1852\,\mathrm m$，一節為一海里每小時，約 $0.51444\,\mathrm{m/s}$。
\ei
\key{閱讀重點}所有計算先確認距離、速度及高度的單位，避免把海里誤寫成速度。
}
\handoutfigure{4}{航電用途與工程單位}{handout02-p004.png}
\slide{自我檢查：航電用途與工程單位}{依 Handout 範圍自編問題與參考解答}{
\small
\key{先自行作答}
自編算例：$180\,\mathrm{kt}$ 等於多少 $\mathrm{km/h}$ 及 $\mathrm{m/s}$？
\key{參考解答}
$180\times1.852=333.36\,\mathrm{km/h}$，再除以 $3.6$ 得 $92.6\,\mathrm{m/s}$。海里是距離，節才是速度。
}
\slide{外界目視與儀表資訊}{Handout 2，PDF pp. 7--10；概念整理與補充推導}{
\small
\bi
\item 低能見度與夜間飛行使姿態、空速、高度及導航資訊更重要。
\item 儀表顯示來源可能是壓力、慣性感測或無線電，不同來源有不同誤差與限制。
\item VFR、IFR 是飛行規則與運作脈絡，不能僅由某一個儀表的存在判定適用資格。
\ei
\key{閱讀重點}儀表協助補足外界線索，但飛行員仍須辨識量測基準及交叉核對。
}
\handoutfigure{9}{外界目視與儀表資訊}{handout02-p009.png}
\slide{自我檢查：外界目視與儀表資訊}{依 Handout 範圍自編問題與參考解答}{
\small
\key{先自行作答}
為什麼姿態指示器與氣壓高度表不能互相取代？
\key{參考解答}
姿態描述機體方向，高度表由壓力推得高度。抬頭姿態不保證爬升，維持同一高度也可能有不同俯仰姿態，兩者提供不同狀態資訊。
}
\section{獨立類比架構}
\slide{獨立儀表與類比傳輸}{Handout 2，PDF pp. 11--17；概念整理與補充推導}{
\small
\bi
\item 傳統獨立儀表通常各有感測器、傳動／傳輸及顯示元件，功能鏈分散。
\item 類比線路可能受到阻抗、雜訊、接地與 EMI 影響，線路與接頭數也增加重量及維修點。
\item 講義的 synchro 案例呈現類比角度傳輸，但個別距離及接線描述不應推成所有類比系統限制。
\ei
\key{閱讀重點}獨立架構容易看清單一功能邊界，整合則要另行設計共享與隔離。
}
\handoutfigure{12}{獨立儀表與類比傳輸}{handout02-p012.png}
\handoutfigure{16}{獨立儀表與類比傳輸}{handout02-p016.png}
\slide{自我檢查：獨立儀表與類比傳輸}{依 Handout 範圍自編問題與參考解答}{
\small
\key{先自行作答}
獨立感測器較多，是否一定比較不可靠？
\key{參考解答}
不能只靠數量判斷。元件多會增加故障點，但獨立來源也可能避免共同依賴。應比較任務可靠度、故障隔離與總體設計。
}
\section{駕駛艙與分工}
\slide{類比與數位顯示的操作差異}{Handout 2，PDF pp. 18--20；概念整理與補充推導}{
\small
\bi
\item 類比儀表將許多狀態分散呈現，飛行員需要掃描與整合各個讀值。
\item 整合顯示可依任務重排資訊，但也增加頁面、模式及來源選擇的辨識需求。
\item 講義用機型與機組分工說明演進，特定年代配置不代表整個系列均相同。
\ei
\key{閱讀重點}比較駕駛艙時，除了畫面數量，也要追蹤資料來源與操作工作量。
}
\handoutfigure{18}{類比與數位顯示的操作差異}{handout02-p018.png}
\handoutfigure{19}{類比與數位顯示的操作差異}{handout02-p019.png}
\slide{自我檢查：類比與數位顯示的操作差異}{依 Handout 範圍自編問題與參考解答}{
\small
\key{先自行作答}
螢幕越大、資訊越多，是否保證工作負荷越小？
\key{參考解答}
不保證。資訊選擇、告警優先序、模式與畫面配置會影響理解速度。資料太密或切換過多都可能增加負擔。
}
\section{資料匯流排}
\slide{數位整合與共享資料}{Handout 2，PDF pp. 21--26；概念整理與補充推導}{
\small
\bi
\item 感測資料轉成數位訊息後，可透過資料匯流排供控制器與顯示器使用。
\item 共享量測可減少重複設備與配線，但要處理格式、時序、更新率與來源一致性。
\item 匯流排協定各有拓樸與方向性，不能把講義的雙向概述套用到每一條航電匯流排。
\ei
\key{閱讀重點}相同的數值若時間不同、單位不同或狀態無效，不能直接當成相同資訊。
}
\handoutfigure{23}{數位整合與共享資料}{handout02-p023.png}
\handoutfigure{25}{數位整合與共享資料}{handout02-p025.png}
\slide{自我檢查：數位整合與共享資料}{依 Handout 範圍自編問題與參考解答}{
\small
\key{先自行作答}
自編算例：每秒發送 100 筆，每筆 64 bytes，純資料率是多少？
\key{參考解答}
$100\times64\times8=51{,}200\,\mathrm{bit/s}$。真正匯流排容量還要容納標頭、校驗、間隔與排程，不能只用純資料率判定足夠。
}
\slide{控制層級與更新率}{Handout 2，PDF pp. 24--28；概念整理與補充推導}{
\small
\bi
\item 飛機穩定控制關注較快的運動，導航與航路管理的變化通常較慢。
\item CPU 時脈、程式執行週期、感測取樣率及控制更新率是不同量。
\item 共享網路還需要限制延遲與資料年齡，平均資料率低仍可能在關鍵時刻延遲。
\item 講義中的組員工時表屬原始版本資料，本自讀筆記不將它當作現行派遣規則。
\ei
\key{閱讀重點}時間需求來自被控制的動態與任務，不能只比較電腦規格。
}
\slide{自我檢查：控制層級與更新率}{依 Handout 範圍自編問題與參考解答}{
\small
\key{先自行作答}
為什麼航路每秒更新一次，不能推論角速度控制也適合相同週期？
\key{參考解答}
航路變化慢，角速度控制需要及時抵抗擾動。過慢的回授會增加相位延遲，降低阻尼與穩定裕度。實際更新率仍須依設計模型決定。
}
\section{系統功能}
\slide{性能、可用性與機地分工}{Handout 2，PDF pp. 29--36；概念整理與補充推導}{
\small
\bi
\item 航電透過導航、控制與系統管理支援飛行任務，效益須連到具體操作需求。
\item 更長航程或更低工作負荷通常還依賴氣動、發動機、程序與人員訓練。
\item 機載系統關注本機操作，地面系統關注空域交通，兩者交換資訊但目的與責任不同。
\ei
\key{閱讀重點}評估整合成效時，可比較任務完成、維修停機及人員操作需求。
}
\slide{自我檢查：性能、可用性與機地分工}{依 Handout 範圍自編問題與參考解答}{
\small
\key{先自行作答}
為什麼「裝數位航電就會增加航程」不是完整工程論證？
\key{參考解答}
需要指出機制，例如減重、改善航路或燃油管理，並在相同任務與條件下比較。航程也受到載重、氣象與推進效率影響。
}
\section{CNS/ATM}
\slide{傳統與衛星資料鏈 CNS}{Handout 2，PDF pp. 37--41；概念整理與補充推導}{
\small
\bi
\item 傳統通訊依賴語音，導航常用地面台，監視常用雷達，講義將其與衛星和資料鏈方案比較。
\item GNSS 提供定位基礎，增強系統改善特定服務能力，ADS-B 廣播機上狀態。
\item 新舊設施可並存，是否能執行某項運作仍取決於整套設備、服務與程序。
\ei
\key{閱讀重點}通訊、導航與監視的功能不同，即使共用衛星或數位技術也須分開分析。
}
\handoutfigure{39}{傳統與衛星資料鏈 CNS}{handout02-p039.png}
\handoutfigure{40}{傳統與衛星資料鏈 CNS}{handout02-p040.png}
\slide{自我檢查：傳統與衛星資料鏈 CNS}{依 Handout 範圍自編問題與參考解答}{
\small
\key{先自行作答}
ADS-B 與雷達得到位置資訊的方式有何主要差別？
\key{參考解答}
ADS-B 使用航空器廣播的資料，依賴機上來源。雷達依回波或合作式應答取得監視資訊。兩者的量測方式、依賴與誤差來源不同。
}
\section{模組化與顯示}
\slide{IMA、LRU 與微電子技術}{Handout 2，PDF pp. 42--44；概念整理與補充推導}{
\small
\bi
\item 講義以模組化、VLSI、MEMS 與處理器解釋航電小型化與資料共享。
\item LRU 是航線維修可更換單元的概念。IMA 關注在共用計算資源上整合功能與隔離。
\item 功能獨立、實體可更換及計算分割是不同層次，不能直接把每個 IMA 功能等同一個 LRU。
\ei
\key{閱讀重點}共享資源需要故障隔離與資源控制，避免單一功能拖累其他功能。
}
\handoutfigure{44}{IMA、LRU 與微電子技術}{handout02-p044.png}
\slide{自我檢查：IMA、LRU 與微電子技術}{依 Handout 範圍自編問題與參考解答}{
\small
\key{先自行作答}
一台共用電腦執行多項功能時，軟體分割能否取代獨立電源？
\key{參考解答}
不能。軟體分割可限制程式之間的影響，但電源失效仍可能讓整台電腦停止。隔離與備援需要涵蓋不同故障層次。
}
\slide{不同機型的座艙閱讀方法}{Handout 2，PDF pp. 45--53；概念整理與補充推導}{
\small
\bi
\item 先辨認姿態、速度、高度、導航、發動機與告警資訊的所在位置，再比較顯示整合程度。
\item 讀取某個數值時，確認量測來源、單位、模式與有效性，不能只認畫面位置。
\item 機型照片可用於架構比較，不能單憑照片推論全部控制律或備援能力。
\ei
\key{閱讀重點}有效的比較以功能、資料來源及飛行階段為軸，不以螢幕數量排名。
}
\handoutfigure{46}{不同機型的座艙閱讀方法}{handout02-p046.png}
\handoutfigure{48}{不同機型的座艙閱讀方法}{handout02-p048.png}
\handoutfigure{51}{不同機型的座艙閱讀方法}{handout02-p051.png}
\slide{自我檢查：不同機型的座艙閱讀方法}{依 Handout 範圍自編問題與參考解答}{
\small
\key{先自行作答}
選兩張座艙圖，列出同一飛行任務需要的資訊，並說明如何交叉核對。
\key{參考解答}
例如進場時核對空速、高度、姿態、橫向與垂直偏差及發動機狀態。交叉核對要辨別獨立來源，共用來源的兩個顯示不構成獨立驗證。
}
\section{維修與 BIT}
\slide{間歇故障、紀錄與更換}{Handout 2，PDF pp. 54--56；概念整理與補充推導}{
\small
\bi
\item BIT 在系統中進行測試與監視，可記錄飛行中發生的故障及事件。
\item 維修人員下載紀錄、辨認時間與條件、隔離可疑單元，再依程序更換及驗證。
\item 故障碼是證據，單元回報異常也可能來自上游感測器、電源或接線問題。
\ei
\key{閱讀重點}維修應追蹤事件脈絡，避免只看單一故障碼就反覆更換正常設備。
}
\handoutfigure{54}{間歇故障、紀錄與更換}{handout02-p054.png}
\slide{自我檢查：間歇故障、紀錄與更換}{依 Handout 範圍自編問題與參考解答}{
\small
\key{先自行作答}
故障只在飛行中出現，落地後功能正常，BIT 可以提供什麼？
\key{參考解答}
事件時間、工作模式、量測有效性及相關故障碼可縮小條件。再比對環境、連接與電源等因素，依維修程序排查，最後確認更換或修復結果。
}
\section{整合複習}
\slide{端到端航電分析}{Handout 2，PDF pp. 57；概念整理與補充推導}{
\small
\bi
\item 以一項任務追蹤感測、處理、資料傳輸、顯示、控制與維修紀錄。
\item 比較整合前後的重量、配線、共享能力與故障傳播途徑。
\item 講義最後把視角延伸到地面 ATC，提醒機上設備屬於更大的資訊系統。
\ei
\key{閱讀重點}完整的答案要同時包含整合效益與隔離、備援及人因需求。
}
\handoutfigure{57}{端到端航電分析}{handout02-p057.png}
\slide{自我檢查：端到端航電分析}{依 Handout 範圍自編問題與參考解答}{
\small
\key{先自行作答}
設計共用空速來源時，至少列出四項應檢查的事項。
\key{參考解答}
可檢查感測器備援與交叉比較、資料單位與時間戳、有效性傳播、電源及匯流排隔離，以及失效後顯示與控制律的降級行為。
}
\section{複習與來源}
\slide{整合複述檢查}{Handout 2 全份複習}{
\bi
\item 選一項系統功能，說明輸入、輸出、座標、單位與資料來源。
\item 找出一個公式，逐一解釋每個符號，再說明它使用哪些假設。
\item 選一個故障或誤差，追蹤它如何影響下游功能，以及可以用什麼獨立資訊核對。
\item 重做本份的自我檢查，再對照講義與相關公式。
\ei
}
\slide{來源與後續閱讀}{Handout 2 與隨附教材}{
\small
\bi
\item 主要來源：賴盈誌老師，CA7528 Digital Avionics，Handout 2。原始檔保留於上一層資料夾。
\item 版型與重用內容：既有 Chapter03\_SelfStudy\_LaTeX。Handout 3、4 沿用相關章節及其原圖。
\item 第三章重用內容來源：R. P. G. Collinson，\textit{Introduction to Avionics Systems}，第 3 版。
\item 本筆記的重算、核對與延伸為自讀整理，並非授課教師或出版社正式勘誤。
\item 原圖內來源標示與圖說保留，新增題目與參考解答已另行標明。
\ei
}
\end{document}
